% ---------
% --------
%!TEX TS-program = pdflatex
%!TEX encoding = UTF-8 Unicode
%!TEX spellcheck = English (Aspell)
\documentclass[11pt]{article}
\usepackage{lucy,handout,graphicx}
\usepackage[utf8]{inputenc}		% Allow some non-ASCII Unicode in source
\usepackage{tikz}

\def\Cdot{\mathbin{\text{\normalfont \textbullet}}}
\def\Sym#1{\texttt{\color{BrickRed}#1}}
\def\BlueSym#1{\textbf{\texttt{\color{RoyalBlue}#1}}}
\allowdisplaybreaks

\newtheorem{theorem}{Claim}

% =========================================================
%   Define common stuff for solution headers
% =========================================================
\Class{CSCI 332}
\Semester{Fall 2026}

% =========================================================
\begin{document}

% ---------------------------------------------------------


\headers{CSCI 332}{Homework 1 Solutions}{Fall 2025}

\begin{enumerate}
\item
For your upcoming family reunion, you are in charge of one of the day's activities.
You decide to have the family head out in two-person canoes. However, you know that
you need to be careful in pairing your family members up, because they can get jealous
if they get stuck spending time with one person when they would prefer to spend time with another.
Because you are taking CSCI 332, you wonder if you could re-work the stable matching
definition from the hospital-student problem to work for your canoe pair problem. You consider
asking each family member to order the rest of the family members according to their
preference for a canoe partner.

For example, for $n=4$ family reunion attendees, you might get the following list:

\begin{center}
\begin{tabular}{ c c c }
    Alice: Ben, Cleo, Desiree \\
    Ben: Cleo, Alice, Desiree \\
    Cleo: Alice, Ben, Desiree\\
    Desiree: Alice, Ben, Cleo
\end{tabular}
\end{center}


\begin{enumerate}[(a)]
    \item 
Re-work the definition of stable matching from the hospital-medical student
        version of the problem to the canoe pair version of the problem.
        You should describe the input of the problem and the desired output of the problem.

\begin{Solution}
    The canoe pair stable matching problem is defined as follows:

    \noindent \emph{Input:} A set of $n$ family members, each of whom ranks the other
    $n-1$ players in order of preference to be their canoe partner. Assume that $n$
    is even.

    \noindent \emph{Output:} A matching of the $n$ family members into $n/2$
    pairs that is \emph{perfect} (every family member is in exactly one pair) and \emph{stable}: that is,
    no two family members who are not currently paired up prefer one another over their current
    pair.
    
\end{Solution}

\item 

Does a stable matching always exist for the canoe pair problem?


\begin{Solution}

    No.
    
\end{Solution}

\item 
How do you know?

\begin{Solution}

   To show that there is not always a stable matching, we can give an example
   input where no such matching exists. Conveninently, the example above is such an input!
   To see why there is no stable matching, consider every perfect matching possible
   for this input.
   \begin{itemize}
    \item $M = \{\{A,B\}, \{C, D\}\}$: in this case, $B$ and $C$ form an
    unstable pair, since $B$ prefers $C$ over $A$ and $C$ prefers $B$ over $D$.
    \item $M = \{\{A,C\}, \{B, D\}\}$: in this case, $A$ and $B$ form an
    unstable pair, since $A$ prefers $B$ over $C$ and $B$ prefers $A$ over $D$.
    \item $M = \{\{A,D\}, \{B, C\}\}$: in this case, $A$ and $C$ form an
    unstable pair, since $A$ prefers $C$ over $D$ and $C$ prefers $A$ over $B$.
   \end{itemize} 
\end{Solution}


\end{enumerate}
\item Back to the hospital-student stable matching problem: describe
and analyze an efficient algorithm to determine whether a given set of preferences
has a \emph{unique} stable matching.

\begin{Solution}
As we saw in class, the Gale-Shapley algorithm produces the ``best'' matching for the proposers
and the ``worst'' matching for the proposees. So we can run Gale-Shapley twice, once with hosptials proposing
and once with students proposing. If the resulting matchings are the same, then it is the unique stable matching
for the input. Otherwise, there are at least two distinct matchings. 

Each run of Gale-Shapley takes $O(n^2)$ time, and checking whether the outputs are equal is also $O(n^2)$ time,
so the runtime is $O(n^2)$ overall.
\end{Solution}

\end{enumerate}

\end{document}

